\chapter{\texorpdfstring{The target language \lbox\ and its semantics}{The target language lambda-box and its semantics}}
\label{chap:lambdabox}

\section*{At a glance}
\begin{itemize}
\item \lead{Goal} Fix \lbox{}'s syntax and flag-parametric, weak call-by-value semantics, independently of Lean-the-source-language.
\item \lead{Main results} Theorem~\ref{thm:value_final} (normal forms evaluate to themselves); Theorem~\ref{thm:eval_to_value} (evaluation yields a value); Theorem~\ref{thm:eval_deterministic} (at most one value); Theorem~\ref{thm:WcbvEval-lbClosed} (evaluation preserves closedness); Theorem~\ref{thm:lbEval_sound} (the executable evaluator is sound).
\item \lead{Status} \stProved{} throughout: nothing here mentions \code{Lean.Expr}, \TrExprS{} or \code{HasType}, so nothing inherits lean4lean's \code{sorryAx} or a trust edge to Chapter~\ref{chap:trust}.
\item \lead{Lean files} \code{Basic.lean}, \code{Printing.lean}, \code{Closed.lean}, and under \code{Semantics/}: \code{Flags.lean}, \code{Env.lean}, \code{Values.lean}, \code{Eval.lean}, \code{Substitution.lean}, \code{Metatheory.lean}, \code{Compute.lean}.
\item \lead{Paper} [JACM, \S7.1, \S7.4] for \WcbvEval{}; Letouzey's amendments [L, Def.~5, Def.~8] for the box and singleton-case rules.
\end{itemize}

\section{\texorpdfstring{The syntax of \lbox}{The syntax of lambda-box}}
\label{sec:lambdabox:syntax}

\begin{definition}[Kernames and module paths]
\label{def:Kername}
\lean{Ident, DirPath, ModPath, Kername, toModPath, cleanIdent, toKername, rootKername}
\leanok
\lead{In short} A \lbox{} \hl{kername} pairs a module path with a base identifier, mirroring MetaRocq's \code{kername}.
\begin{itemize}
\item module path: \code{MPfile dp} or \code{MPdot mp id} (\code{MPbound}, for ML functor arguments, is absent, Lean having none);
\item \code{toKername}: path from the prefix, identifier from the last component, \code{cleanIdent}-escaped; \code{rootKername s}: unqualified, identifier \code{cleanIdent s}.
\end{itemize}
\end{definition}

\begin{remark}
\lead{Caveat} \code{toKername} is \textbf{not injective}: \code{cleanIdent} maps an already-escaped \texttt{\_u}\ldots string onto what it escapes, and distinct \code{.num} components can collapse to one decimal spelling (F-KERNAME, \code{toKername\_not\_injective}). It panics on the anonymous name.
\end{remark}

\begin{definition}[Binder annotations]
\label{def:BinderName}
\lean{BinderName}
\leanok
\lead{In short} MetaRocq's \code{name}: a bound variable is either \hl{printable} (\code{named s}) or \code{anon}; inert for the semantics (no rule inspects it), essential for the printer, which rejects names outside printable ASCII (a well-formedness clause in Chapter~\ref{chap:lowering}).
\end{definition}

\begin{definition}[Inductive identifiers and projections]
\label{def:InductiveId}
\lean{InductiveId, ProjectionInfo}
\leanok
\uses{def:Kername}
\lead{In short} \code{InductiveId}: a mutual block's kername plus the member's index (MetaRocq's \code{inductive} pair). \code{ProjectionInfo}: the inductive projected from, its parameter count, and the field's index, counted \hl{after} pruning.
\end{definition}

\begin{definition}[Members of a mutual fixpoint block]
\label{def:FixDef}
\lean{FixDef}
\leanok
\uses{def:BinderName}
\lead{In short} One member of a mutual fixpoint block: a display name, a body, and guard index \code{principalArgIdx} (default $0$); the body must begin with at least $\mathtt{principalArgIdx}+1$ lambdas.
\lead{Lean} \hl{The guard index} is what the guarded-fixpoint rules of Section~\ref{sec:lambdabox:eval} count arguments against.
\end{definition}

\begin{remark}
\lead{Caveat} The source calls \code{principalArgIdx} irrelevant, safe at $0$: true for the \emph{printer}, false for the guarded-fix rules, where it decides when unfolding fires. The eraser always leaves it at $0$ (F-ETA, Chapter~\ref{chap:eraser}).
\end{remark}

\begin{definition}[Primitive values]
\label{def:PrimVal}
\lean{PrimTag, PrimModel, PrimVal}
\leanok
\lead{In short} The only modelled primitive tag is \code{primInt} ($\mathtt{BitVec}\ 63$), a dependent pair of tag and model element. MetaRocq's \code{primFloat}/\code{primArray} are commented out (arrays need a positivity trick avoided here).
\end{definition}

\begin{definition}[\texorpdfstring{The \lbox{} term language}{The lambda-box term language}]
\label{def:LBTerm}
\lean{LBTerm}
\leanok
\uses{def:BinderName, def:Kername, def:InductiveId, def:FixDef, def:PrimVal}
\lead{In short} \hl{\code{LBTerm}} is MetaRocq's \code{EAst.term}, twelve constructors, formalised independently of Lean's source syntax.
\begin{center}
\begin{tabular}{lp{0.66\linewidth}l}
\toprule
\textbf{constructor} & \textbf{meaning} & \textbf{MetaRocq} \\
\midrule
\code{box} & the box $\Box$, erased-proof placeholder & \code{tBox} \\
\code{bvar} & de Bruijn variable & \code{tRel} \\
\code{fvar} & free variable, carries a \code{Lean.FVarId} & --- (added) \\
\code{lambda} & abstraction, binder annotation & \code{tLambda} \\
\code{letIn} & let binding, binder annotation & \code{tLetIn} \\
\code{app} & binary application & \code{tApp} \\
\code{const} & constant, keyed by kername & \code{tConst} \\
\code{construct} & inductive id, index, \emph{list} of arguments & \code{tConstruct} \\
\code{case} & inductive+param.\ count, discriminee, branches & \code{tCase} \\
\code{proj} & primitive projection & \code{tProj} \\
\code{fix} & mutual fixpoint block, selected member & \code{tFix} \\
\code{prim} & primitive value & \code{tPrim} \\
\bottomrule
\end{tabular}
\end{center}
\lead{Lean} No \code{cofix}: Lean has no coinductive types, so MetaRocq's cofixpoint rules are unrepresentable here.
\end{definition}

\begin{remark}
\lead{Deviation} Against \code{E.term} of [JACM, Fig.~16]: \code{.fvar} is added (locally nameless, [R, \S4.7]), branches carry binder-name lists ([JACM, Fig.~18]), and \code{tCoFix} is dropped (\code{.prim} mirrors MetaRocq's own later addition, not a fourth change). Block vs.\ applied form is a \emph{flag} (Section~\ref{sec:lambdabox:flags}); the eraser emits applied form.
\end{remark}

\begin{definition}[Abstraction of a free variable]
\label{def:toBvar}
\lean{toBvar, toBvarArgs, toBvarAlts, toBvarDefs, abstract}
\leanok
\uses{def:LBTerm}
\lead{In short} $\mathtt{toBvar}\ x\ lvl\ t$ replaces free variable $x$ in $t$ by de Bruijn index $lvl$; $\mathtt{abstract}\ x\ e = \mathtt{toBvar}\ x\ 0\ e$ closes a binder the eraser just emitted.
\begin{itemize}
\item $lvl$ rises by one under a \code{lambda}/\code{letIn} body, by a branch's field count under a case alternative, by the block size under a \code{fix};
\item the three list companions are written by hand, not as \code{List.map}, so the structural-recursion checker sees through the nested lists.
\end{itemize}
\end{definition}

\subsection*{The printer, and where the verified perimeter stops}

\code{Printing.lean}'s S-expression printer is \hl{outside the verified perimeter}: \code{to\_sexpr} is a \code{partial def} with no equational lemmas, mentioned nowhere else in the development, and its \code{.fvar} case is \code{unreachable!}. Every claim below of the form ``the emitted program satisfies \ldots'' concerns the \code{Program} \emph{value} of Definition~\ref{def:ASTType}, not parsed bytes.

\section{Declarations, environments and programs}
\label{sec:lambdabox:env}

\begin{definition}[Inductive declarations]
\label{def:OneInductiveBody}
\lean{ConstructorBody, ProjectionBody, AllowedEliminations, OneInductiveBody, RecursivityKind, MutualInductiveBody}
\leanok
\uses{def:Kername}
\lead{In short} What \lbox{} records about an inductive: a \code{ConstructorBody} (display name, field count \emph{after} pruning --- \code{nargs}, not the arity of Definition~\ref{def:constructorArity}, which adds parameter count), a \code{OneInductiveBody} (name, \hl{\code{propositional} mark}, \code{kelim}, constructors, projections), a \code{MutualInductiveBody} (recursivity kind, \code{npars}, bodies).
\begin{itemize}
\item \code{propositional} has \emph{no default}: the \code{iota\_sing}/\code{proj\_prop} rules (Section~\ref{sec:lambdabox:eval}) read it to collapse an elimination of a boxed discriminee, so every producer must supply the right value;
\item \code{kelim} and the recursivity kind are carried for fidelity to MetaRocq's format; no pass in MetaRocq's erasure pipeline or in peregrine reads either to decide anything.
\end{itemize}
\end{definition}

\begin{definition}[\texorpdfstring{The \lbox{} global environment}{The lambda-box global environment}]
\label{def:GlobalDeclarations}
\lean{ConstantBody, GlobalDecl, GlobalDeclarations}
\leanok
\uses{def:LBTerm, def:OneInductiveBody, def:Kername}
\lead{In short} MetaRocq's \code{global\_declarations}: kername/declaration pairs, each a constant (\emph{optional} body; \code{none} is an unrealisable axiom) or a mutual inductive block, in \hl{reverse} order (first-added deepest), so head-first lookup finds the most recent binding.
\end{definition}

\begin{definition}[Programs]
\label{def:ASTType}
\lean{ASTType, Program}
\leanok
\uses{def:GlobalDeclarations, def:LBTerm}
\lead{In short} A \lbox{} program on disk: an environment plus an optional main term; \code{untyped} restricts peregrine's \code{PAst} to \code{Untyped}, so this frontend never emits typed \lboxT{}. \code{Program} (a synonym) is the eraser's return type, the capstone's subject (Chapter~\ref{chap:capstone}).
\end{definition}

\section{The de Bruijn kit}
\label{sec:lambdabox:subst}

This section's declarations live in \code{Semantics/Substitution.lean}: \code{envLookup}, \code{shift} and companions are root names; \code{mkLambdas} sits in \code{namespace LeanToLambdaBox}.

\begin{definition}[Environment lookup]
\label{def:LBTerm-envLookup}
\lean{ModPath.beq, Kername.beq, LBTerm.envLookup}
\leanok
\uses{def:GlobalDeclarations, def:Kername}
\lead{In short} $\mathtt{envLookup}\ \Gamma\ kn$ returns the \hl{first} entry's declaration with kername \code{Kername.beq}-equal to $kn$: MetaRocq's \code{lookup\_env}. Delta reduction, the propositional mark, and constructor arities all read the environment through it.
\end{definition}

\begin{definition}[Index shifting]
\label{def:LBTerm-shift}
\lean{LBTerm.shift, LBTerm.shiftArgs, LBTerm.shiftAlts, LBTerm.shiftDefs}
\leanok
\uses{def:LBTerm}
\lead{In short} $\mathtt{shift}\ d\ c\ t$ raises every index of $t$ at or above cutoff $c$ by $d$, pushing $c$ up by one under a \code{lambda}/\code{letIn} body, by a branch's field count under an alternative, by the block size under a \code{fix}.
\lead{Lean} Pinned to lean4lean's loose-index lift, lining up source/target de Bruijn bookkeeping in the erasure relation (Chapter~\ref{chap:erases}).
\end{definition}

\begin{definition}[Substitution]
\label{def:LBTerm-subst}
\lean{LBTerm.subst, LBTerm.substArgs, LBTerm.substAlts, LBTerm.substDefs, LBTerm.subst1, LBTerm.substList}
\leanok
\uses{def:LBTerm-shift}
\lead{In short} $\mathtt{subst}\ s\ d\ t$ replaces the bound variable at depth $d$ by $s$ (shifted by $d$), decrementing every higher index; pinned to lean4lean's single-variable instantiation.
\begin{itemize}
\item $\mathtt{subst1}\ s\ t$: depth-$0$ case, the redex of beta and zeta;
\item $\mathtt{substList}\ ss\ t$: folds \code{subst1} left to right, MetaRocq's \code{substl}, the redex of the iota rules and fixpoint unfolding.
\end{itemize}
\end{definition}

\begin{definition}[Application spines]
\label{def:LBTerm-mkApps}
\lean{LBTerm.mkApps, LBTerm.spineHead, LBTerm.spineArgs}
\leanok
\uses{def:LBTerm}
\lead{In short} $\mathtt{mkApps}\ f\ [a_1,\dots,a_n]$ is the left-nested spine $(\dots((f\,a_1)\,a_2)\dots a_n)$, MetaRocq's \code{mkApps}; \code{spineHead} peels every application node (\code{EAstUtils.head}), \code{spineArgs} collects the arguments in order.
\lead{Lean} The applied-form constructor and fixpoint-spine values of Section~\ref{sec:lambdabox:values} are exactly \code{mkApps}-shapes.
\end{definition}

\begin{lemma}[Spine decomposition]
\label{lem:LBTerm-mkApps_concat}
\lean{LBTerm.mkApps_concat, LBTerm.spineHead_mkApps, LBTerm.spineArgs_mkApps}
\leanok
\uses{def:LBTerm-mkApps}
\lead{In short} Extending a spine by one argument:
\begin{itemize}
\item $\mathtt{mkApps}\ f\ (args \mathbin{+\mkern-10mu+} [a]) = \mathtt{app}\ (\mathtt{mkApps}\ f\ args)\ a$;
\item $\mathtt{spineHead}\ (\mathtt{mkApps}\ f\ args) = \mathtt{spineHead}\ f$;
\item $\mathtt{spineArgs}\ (\mathtt{mkApps}\ f\ args) = \mathtt{spineArgs}\ f \mathbin{+\mkern-10mu+} args$.
\end{itemize}
\end{lemma}
\begin{proof}
\leanok
\uses{def:LBTerm-mkApps}
\lead{Idea} Induction on the argument list, unfolding \code{mkApps}/\code{spineHead}/\code{spineArgs} at \code{.app}; \code{simp} closes the append cases.
\end{proof}

\begin{lemma}[Spine injectivity and disjointness]
\label{lem:LBTerm-mkApps_construct_inj}
\lean{LBTerm.mkApps_construct_inj, LBTerm.mkApps_fix_inj, LBTerm.mkApps_construct_ne_fix}
\leanok
\uses{def:LBTerm-mkApps}
\lead{In short} An applied-form spine determines its data: if $\mathtt{mkApps}\ (\mathtt{construct}\ iid\ c\ [])\ args = \mathtt{mkApps}\ (\mathtt{construct}\ iid'\ c'\ [])\ args'$ then $iid=iid'$, $c=c'$, $args=args'$; likewise a fixpoint spine determines block, member, arguments; a constructor spine is never a fixpoint spine.
\lead{Lean} Replaces case analysis on an \code{mkApps}-indexed relation, used by the determinism proof (Theorem~\ref{thm:eval_deterministic}).
\end{lemma}
\begin{proof}
\leanok
\uses{lem:LBTerm-mkApps_concat}
\lead{Idea} Apply the spine-head/argument lemmas to both sides; the heads reduce to the bare \code{construct}/\code{fix} node by \code{simp}, giving the data by injectivity. Disjointness: distinct constructors of \code{LBTerm}.
\end{proof}

\begin{definition}[The fixpoint-unfolding substitution]
\label{def:LBTerm-fixSubst}
\lean{LBTerm.fixSubst}
\leanok
\uses{def:LBTerm, def:FixDef}
\lead{In short} $\mathtt{fixSubst}\ defs = [\mathtt{fix}\ defs\ (n-1),\dots,\mathtt{fix}\ defs\ 0]$, $n=|defs|$: MetaRocq's \code{fix\_subst}; index $i$ of a body refers to member $n-1-i$ of its block.
\lead{Lean} Invisible for a single-member block, needed for a mutual one, pinned by the two-member witness of Lemma~\ref{lem:mf_fixSubst}.
\end{definition}

\begin{definition}[\texorpdfstring{A \code{Prop}-motive recursor for \code{LBTerm}}{A Prop-motive recursor for LBTerm}]
\label{def:LBTerm-recData}
\lean{LBTerm.recData}
\leanok
\uses{def:LBTerm}
\lead{In short} \code{LBTerm} is a \emph{nested} inductive (lists inside \code{construct}, \code{case}, \code{fix}), so plain induction is rejected; \code{recData} replaces it, an \code{elab\_as\_elim} eliminator with a \code{Prop} motive whose list-carrying arms hand back membership hypotheses instead of raw nested motives.
\lead{Lean} Defined by \code{refine} on the kernel recursor with six motives; used throughout the later chapters.
\end{definition}

\begin{lemma}[The hand-rolled traversals are maps]
\label{lem:LBTerm-shiftArgs_eq_map}
\lean{LBTerm.shiftArgs_eq_map, LBTerm.substArgs_eq_map, LBTerm.shiftAlts_eq_map, LBTerm.substAlts_eq_map, LBTerm.shiftDefs_eq_map, LBTerm.substDefs_eq_map}
\leanok
\uses{def:LBTerm-shift, def:LBTerm-subst}
\lead{In short} Each of the six list companions of \code{shift}/\code{subst} equals the corresponding \code{List.map}.
\begin{itemize}
\item arguments: pointwise $\mathtt{shift}\ d\ c$, resp.\ $\mathtt{subst}\ s\ d$; case alternatives: binder list kept, cutoff/depth raised by the branch's field count;
\item fixpoint definitions: only the body traversed, cutoff unmoved (bodies already live under the block's binders).
\end{itemize}
\end{lemma}
\begin{proof}
\leanok
\uses{def:LBTerm-shift, def:LBTerm-subst}
\lead{Idea} List induction, unfolding one step each case; the bridge every \code{recData} induction needs.
\end{proof}

\begin{definition}[Lambda telescopes]
\label{def:mkLambdas}
\lean{LeanToLambdaBox.mkLambdas}
\leanok
\uses{def:LBTerm, def:BinderName}
\lead{In short} $\mathtt{mkLambdas}\ names\ body$ wraps $body$ in one \code{lambda} per binder name, outermost first: a case alternative's $(names,body)$ read back as a minor function, how the lowering pass (Chapter~\ref{chap:lowering}) turns branches into functions.
\end{definition}

\begin{lemma}[Telescopes commute with the de Bruijn kit]
\label{lem:shift_mkLambdas}
\lean{LeanToLambdaBox.shift_mkLambdas, LeanToLambdaBox.subst_mkLambdas}
\leanok
\uses{def:mkLambdas, def:LBTerm-shift, def:LBTerm-subst}
\lead{In short} Shifting/substitution pass through a lambda telescope, raising the cutoff (resp.\ depth) by the telescope's length:
$$\mathtt{shift}\ d\ c\ (\mathtt{mkLambdas}\ names\ body) = \mathtt{mkLambdas}\ names\ (\mathtt{shift}\ d\ (c+|names|)\ body)$$
and likewise for \code{subst} at depth $d + |names|$.
\end{lemma}
\begin{proof}
\leanok
\uses{def:mkLambdas, def:LBTerm-shift, def:LBTerm-subst}
\lead{Idea} Induction on the list of names, unfolding one \code{lambda} per step and re-associating $+1$ into the telescope length.
\end{proof}

\subsection*{Closedness and the commutation kit}

The remaining declarations live in \code{LeanToLambdaBox/Closed.lean}, independent of the erasure relation.

\begin{definition}[De Bruijn closedness of a target term]
\label{def:LBClosed}
\lean{LeanToLambdaBox.LBClosed, LeanToLambdaBox.LBClosedArgs, LeanToLambdaBox.LBClosedAlts, LeanToLambdaBox.LBClosedDefs, LeanToLambdaBox.LBClosedArgs_iff, LeanToLambdaBox.LBClosedAlts_iff, LeanToLambdaBox.LBClosedDefs_iff}
\leanok
\uses{def:LBTerm}
\lead{In short} $\mathtt{LBClosed}\ t\ k$ holds when no loose index $\ge k$ occurs in $t$, matching lean4lean's closedness predicate by the same recursion as \code{shift}. Free variables are \hl{not} constrained (a different predicate's job).
\end{definition}

\begin{lemma}[Shift and substitution are inert on closed terms]
\label{lem:LBClosed-shift_eq}
\lean{LeanToLambdaBox.LBClosed.shift_eq, LeanToLambdaBox.LBClosed.subst_eq, LeanToLambdaBox.lbClosed_of_shift_eq}
\leanok
\uses{def:LBClosed, def:LBTerm-shift, def:LBTerm-subst}
\lead{In short} If $t$ is closed below $k$ and cutoff $c \ge k$, neither $\mathtt{shift}\ d\ c$ nor $\mathtt{subst}\ s\ c$ touches any index of $t$; conversely, at unit shift with the bound as cutoff, a term fixed by the shift has no loose index at or above it --- a shift-inertness \emph{equation} turned into a closedness \emph{fact}.
\end{lemma}
\begin{proof}
\leanok
\uses{def:LBClosed, def:LBTerm-recData, lem:LBTerm-shiftArgs_eq_map}
\lead{Idea} One \code{recData} induction, threading $k \le c$ under each binder (the variable case is arithmetic); the converse reads it backwards.
\end{proof}

\begin{lemma}[The arithmetic of closedness]
\label{lem:LBClosed-substList}
\lean{LeanToLambdaBox.LBClosed.mono, LeanToLambdaBox.LBClosed.shift, LeanToLambdaBox.LBClosed.subst_gen, LeanToLambdaBox.LBClosed.subst, LeanToLambdaBox.LBClosed.subst1, LeanToLambdaBox.LBClosed.substList, LeanToLambdaBox.LBClosed.mkApps, LeanToLambdaBox.LBClosed.mkApps_inv, LeanToLambdaBox.LBClosed.mkLambdas, LeanToLambdaBox.LBClosed.mkLambdas_inv}
\leanok
\uses{def:LBClosed, def:LBTerm-mkApps, def:mkLambdas}
\begin{itemize}
\item closedness is monotone in the bound; a shift by $d$ raises it by $d$;
\item substituting a \hl{closed} term at depth $k$ lowers the bound from $k+1$ to $k$; a list of closed terms closes the term outright;
\item closedness passes through the spine and telescope builders both ways.
\end{itemize}
\end{lemma}
\begin{proof}
\leanok
\uses{def:LBClosed, lem:LBClosed-shift_eq, def:LBTerm-recData}
\lead{Idea} Induction with the bound generalised; the substitution laws case-split on the index/depth comparison, the inertness lemma discharging the miss.
\end{proof}

\begin{lemma}[The de Bruijn commutation kit]
\label{lem:LBTerm-subst_subst}
\lean{LeanToLambdaBox.LBTerm.subst_subst, LeanToLambdaBox.LBTerm.subst_subst_gen, LeanToLambdaBox.LBTerm.shift_shift, LeanToLambdaBox.LBTerm.subst_shift_cancel, LeanToLambdaBox.LBTerm.subst_shift_comm, LeanToLambdaBox.LBTerm.substList_append, LeanToLambdaBox.LBTerm.substList_concat}
\leanok
\uses{def:LBTerm-shift, def:LBTerm-subst}
\begin{itemize}
\item two shifts collapse into one when the outer cutoff lies inside the inner shift's band; a substitution inside that band lowers the shift by one; a substitution commutes with an outer shift, pushing its depth;
\item distribution: $\mathtt{subst}\ s\ d\ (\mathtt{subst}\ t\ 0\ u) = \mathtt{subst}\ (\mathtt{subst}\ s\ d\ t)\ 0\ (\mathtt{subst}\ s\ (d+1)\ u)$;
\item simultaneous substitution splits along list append and a trailing element.
\end{itemize}
\end{lemma}
\begin{proof}
\leanok
\uses{def:LBTerm-recData, lem:LBTerm-shiftArgs_eq_map, def:LBTerm-shift, def:LBTerm-subst}
\lead{Idea} Induction over the target syntax with both depths generalised; arithmetic side conditions are carried as equations, so each law has a general form whose instance at depth $0$ is the headline statement.
\end{proof}

\begin{lemma}[Pushing a substitution out through a block of binders]
\label{lem:LBTerm-substList_reverse_subst}
\lean{LeanToLambdaBox.LBTerm.substList_reverse_subst}
\leanok
\uses{def:LBClosed, def:LBTerm-subst}
\lead{In short} Substituting $f$ under the $|rest|$ binders a list is about to fill equals filling them first and substituting $f$ at depth $0$ --- \hl{provided every term in the list is closed}.
\lead{Caveat} Not slack: a counterexample (discriminee $\mathtt{bvar}\ 0$, constant-function body) is recorded; why MetaRocq's evaluation relation carries closedness side conditions throughout.
\end{lemma}
\begin{proof}
\leanok
\uses{lem:LBTerm-subst_subst, lem:LBClosed-shift_eq}
\lead{Idea} Induction on the list; distribution moves the outer substitution past one element, and that element's closedness makes the inner one inert.
\end{proof}

\begin{lemma}[Closedness of the binder-closing folds]
\label{lem:lbClosed_fix_of_bodies}
\lean{LeanToLambdaBox.lbClosed_fix_of_bodies, LeanToLambdaBox.lbClosed_toBvar, LeanToLambdaBox.lbClosed_foldl_toBvar, LeanToLambdaBox.lbClosed_foldl_zipIdx, LeanToLambdaBox.lbClosed_foldl_zipIdx_map}
\leanok
\uses{def:LBClosed, def:toBvar, def:FixDef}
\lead{In short} Folding abstraction over free variables takes a body closed at $0$ to one closed at the list's length; a block closed at its own size is therefore a closed fixpoint node, at every index.
\lead{Caveat} Holds only for a block whose bodies were closed over the block's own names, as the eraser's definition builder does; an arbitrary block can fail (body $\mathtt{bvar}\ 5$ refutes it).
\end{lemma}
\begin{proof}
\leanok
\uses{def:LBClosed, def:toBvar}
\lead{Idea} The fold lemmas are inductions on the list of names; the fixpoint statement is the definitional unfolding of \code{LBClosed}'s fixpoint arm at bound $0$ plus the length equation.
\end{proof}

\section{Evaluation flags and environment queries}
\label{sec:lambdabox:flags}

Evaluation is flag-parametric, as in MetaRocq: the development names four points of the flag cube, used throughout what follows.

\begin{definition}[Evaluation flags]
\label{def:WcbvFlags}
\lean{LeanToLambdaBox.WcbvFlags}
\leanok
\lead{In short} A triple of Booleans parameterising weak call-by-value evaluation, matching MetaRocq's \code{EWcbvEval.WcbvFlags} [JACM, \S7.4].
\begin{center}
\begin{tabular}{p{0.24\linewidth}p{0.52\linewidth}l}
\toprule
\textbf{flag} & \textbf{meaning} & \textbf{at \code{eraseFlags}} \\
\midrule
\code{with\_prop\_case} & a case/projection on an erased proof reduces by substituting boxes & \code{false} \\
\code{with\_guarded\_fix} & a fixpoint unfolds only once its principal argument arrives (Rocq/Lean-style) & \code{true} \\
\code{with\_constructor\_as\_block} & block (args inside \code{construct}) vs.\ applied (args via \code{app}) & \code{false} \\
\bottomrule
\end{tabular}
\end{center}
\end{definition}

\begin{definition}[The four flag points]
\label{def:eraseFlags}
\lean{LeanToLambdaBox.eraseFlags, LeanToLambdaBox.entryFlags, LeanToLambdaBox.blockFlags, LeanToLambdaBox.propBlockFlags}
\leanok
\uses{def:WcbvFlags}
\lead{In short} Four points of the flag cube.
\begin{center}
\begin{tabular}{llp{0.57\linewidth}}
\toprule
\textbf{point} & \textbf{prop, fix, ctor} & \textbf{role} \\
\midrule
\code{eraseFlags} & off, on, applied & \hl{the target semantics}: correctness, capstone, every rung \\
\code{entryFlags} & on, on, applied & MetaRocq default; peregrine's pipeline input \\
\code{blockFlags} & off, on, block & target of \code{optimize} (Ch.~\ref{chap:lowering}) \\
\code{propBlockFlags} & on, on, block & source of \code{optimize} \\
\bottomrule
\end{tabular}
\end{center}
\lead{Lean} \code{eraseFlags} is MetaRocq's \code{opt\_wcbv\_flags}.
\end{definition}

\begin{definition}[The propositional mark]
\label{def:isPropositionalInductive}
\lean{LeanToLambdaBox.isPropositionalInductive}
\leanok
\uses{def:LBTerm-envLookup, def:OneInductiveBody, def:InductiveId}
\lead{In short} Looks the mutual block up, selects the member, returns its \code{propositional} mark (\code{false} on lookup failure): the Boolean half of MetaRocq's \code{inductive\_isprop\_and\_pars}. Guards six rules: \code{false} for \code{iota}/\code{iota\_block}/\code{proj}/\code{proj\_block}, \code{true} for \code{iota\_sing}/\code{proj\_prop}.
\end{definition}

\begin{definition}[Constructor arity]
\label{def:constructorArity}
\lean{LeanToLambdaBox.constructorArity}
\leanok
\uses{def:LBTerm-envLookup, def:OneInductiveBody}
\lead{In short} The declared arity of a constructor: block parameter count plus its own field count, MetaRocq's $\mathtt{cstr\_arity} = \mathtt{ind\_npars}+\mathtt{cstr\_nargs}$; \code{none} when undeclared. Two roles in applied form: the nullary-constructor rule's declaredness condition, and the saturation bound past which a spine stops accumulating.
\end{definition}

\begin{definition}[Collapsible case nodes]
\label{def:wouldCollapse}
\lean{LeanToLambdaBox.wouldCollapse}
\leanok
\uses{def:isPropositionalInductive}
\lead{In short} True when the inductive is propositional and the branches are a singleton: the test for case nodes the singleton rule fires on and \code{optimize} (Chapter~\ref{chap:lowering}) collapses. Read only by that pass.
\end{definition}

\section{Values}
\label{sec:lambdabox:values}

\begin{definition}[Head atoms and head-shape tests]
\label{def:atomValue}
\lean{LeanToLambdaBox.atomValue, LeanToLambdaBox.isLambda, LeanToLambdaBox.isBox, LeanToLambdaBox.isFix, LeanToLambdaBox.isConstruct, LeanToLambdaBox.isPrim, LeanToLambdaBox.isFixApp, LeanToLambdaBox.isConstructApp, LeanToLambdaBox.isPrimApp}
\leanok
\uses{def:LBTerm, def:LBTerm-mkApps}
\lead{In short} $\mathtt{atomValue}\ t$ holds for \lbox{}'s head atoms (box, lambda, fvar, primitive, bare fixpoint) and nothing else: MetaRocq's \code{atom} minus its flag-dependent nullary-\code{tConstruct} clause (below). A constant is \hl{not} an atom, since it delta reduces.
\lead{Lean} MetaRocq's five head-shape tests and spine variants \code{isFixApp}, \code{isConstructApp}, \code{isPrimApp} accompany it; decidable.
\end{definition}

\begin{definition}[Stuck application heads]
\label{def:isStuckApp}
\lean{LeanToLambdaBox.isStuckApp}
\leanok
\uses{def:atomValue, def:WcbvFlags}
\lead{In short} An evaluated function head is \hl{stuck} when no application rule can fire on it: not a lambda, fixpoint (guarded recursion: fixpoint-spine head), box, constructor-spine head, or primitive-application head.
\lead{Lean} MetaRocq's \code{eval\_app\_cong} side condition, negated, minus \code{isLazyApp} (no \code{tLazy} here).
\end{definition}

\begin{definition}[Weak call-by-value normal forms]
\label{def:Value}
\lean{LeanToLambdaBox.Value, LeanToLambdaBox.Value.atom, LeanToLambdaBox.Value.construct_block, LeanToLambdaBox.Value.construct_nil, LeanToLambdaBox.Value.construct_app_val, LeanToLambdaBox.Value.app_stuck, LeanToLambdaBox.Value.fix_app_val}
\leanok
\uses{def:atomValue, def:isStuckApp, def:constructorArity, def:LBTerm-mkApps, def:WcbvFlags, def:FixDef}
\lead{In short} $\mathtt{Value}\ \Gamma\ fl\ t$ is MetaRocq's \code{EWcbvEval.value}: environment-relative, since an applied-form constructor spine is a value only under the declared arity.
\begin{itemize}
\item \code{atom}: a head atom is a value; \code{construct\_block}: block mode, a constructor node with all-value arguments;
\item \code{construct\_nil}/\code{construct\_app\_val}: applied mode, a declared constructor's bare head, resp.\ such a spine extended by a value argument under the arity;
\item \code{app\_stuck}: a stuck head applied to a value (\code{fvar}-headed spines); \code{fix\_app\_val}: guarded recursion, a fixpoint spine under its principal-argument index, applied to a value.
\end{itemize}
\end{definition}

\begin{remark}
\lead{Why} The two spine clauses are stated \textbf{structurally over \code{.app}}, spine shape carried as a premise equation, so \code{Value} stays invertible by case analysis; MetaRocq indexes directly by $\mkApps f\ args$. Compare [JACM, Fig.~12].
\end{remark}

\section{Weak call-by-value evaluation}
\label{sec:lambdabox:eval}

\begin{definition}[Weak call-by-value evaluation]
\label{def:WcbvEval}
\lean{LeanToLambdaBox.WcbvEval, LeanToLambdaBox.WcbvEval.box, LeanToLambdaBox.WcbvEval.lam, LeanToLambdaBox.WcbvEval.fvar, LeanToLambdaBox.WcbvEval.prim, LeanToLambdaBox.WcbvEval.fix_atom, LeanToLambdaBox.WcbvEval.beta, LeanToLambdaBox.WcbvEval.app_box, LeanToLambdaBox.WcbvEval.zeta, LeanToLambdaBox.WcbvEval.delta, LeanToLambdaBox.WcbvEval.construct, LeanToLambdaBox.WcbvEval.construct_atom, LeanToLambdaBox.WcbvEval.construct_app, LeanToLambdaBox.WcbvEval.iota, LeanToLambdaBox.WcbvEval.iota_block, LeanToLambdaBox.WcbvEval.iota_sing, LeanToLambdaBox.WcbvEval.proj, LeanToLambdaBox.WcbvEval.proj_block, LeanToLambdaBox.WcbvEval.proj_prop, LeanToLambdaBox.WcbvEval.fix_guarded, LeanToLambdaBox.WcbvEval.fix_stuck, LeanToLambdaBox.WcbvEval.fix_unguarded, LeanToLambdaBox.WcbvEval.app_cong}
\leanok
\uses{def:WcbvFlags, def:isPropositionalInductive, def:constructorArity, def:isStuckApp, def:LBTerm-envLookup, def:LBTerm-subst, def:LBTerm-mkApps, def:LBTerm-fixSubst, def:LBTerm, def:GlobalDeclarations}
\lead{In short} The target semantics of the whole verification, $\Gamma \vdash t \Downarrow_{fl} v$: \hl{twenty-two rules}, the Lean counterpart of MetaRocq's \code{EWcbvEval.eval} [JACM, \S7.1].
\begin{center}
\begin{tabular}{p{0.22\linewidth}p{0.31\linewidth}p{0.37\linewidth}}
\toprule
\textbf{rule(s)} & \textbf{redex shape} & \textbf{result / side condition} \\
\midrule
\code{box}, \code{lam}, \code{fvar}, \code{prim}, \code{fix\_atom} & a head atom & itself; no reduction under a binder \\
\code{beta} & $f\Downarrow\mathtt{lambda}\ n\ b$, $a\Downarrow av$ & $\mathtt{subst1}\ av\ b$, evaluated \\
\code{app\_box} & $f\Downarrow\Box$, $a\Downarrow$ (any) & $\Box$ ($\Box\,x\to\Box$ [L, Def.~5]; $a$ still evaluated) \\
\code{zeta} & $\mathtt{letIn}\ n\ v\ b$ & $v$'s value substituted into $b$, evaluated \\
\code{delta} & $\mathtt{const}\ k$, $k$ has a body & that body, evaluated (stuck if $k$ is an axiom) \\
\code{construct} (block) & $\mathtt{construct}\ iid\ k\ args$ & $args$ evaluated pointwise \\
\code{construct\_atom} (applied) & bare declared $\mathtt{construct}\ iid\ c\ []$ & itself \\
\code{construct\_app} (applied) & spine under arity, $+1$ arg & arg accumulated onto the spine \\
\code{iota}, \code{iota\_block} (non-\code{Prop}) & ctor-spine (resp.\ block) discriminee & $\mathtt{substList}\ ((args.\mathtt{drop}\ np).\mathtt{reverse})\ body_k$, evaluated \\
\code{iota\_sing} (prop-case) & singleton case, \code{Prop} inductive, discriminee $\Downarrow \Box$ & $|names|$ boxes into the sole branch, evaluated ([L, Def.~8]) \\
\code{proj}, \code{proj\_block} (non-\code{Prop}) & ctor-$0$ (resp.\ block) spine & field $\mathtt{paramCount}+\mathtt{fieldIdx}$, evaluated \\
\code{proj\_prop} (prop-case) & \code{Prop} discriminee $\Downarrow \Box$ & $\Box$ \\
\code{fix\_guarded} (guarded fix) & fix spine, count $=$ \code{principalArgIdx} & unfold via \code{fixSubst}, applied to $argsv$ then the new arg \\
\code{fix\_stuck} (guarded fix) & fix spine strictly under \code{principalArgIdx} & accumulate (a value) \\
\code{fix\_unguarded} (unguarded fix) & bare fixpoint head, any argument & unfold on every argument \\
\code{app\_cong} (head stuck) & $f\Downarrow$ a stuck head, $a\Downarrow$ & application of the two values \\
\bottomrule
\end{tabular}
\end{center}
\lead{Lean} No \code{cofix} rule: \code{LBTerm} has no cofixpoint node, so MetaRocq's \code{eval\_cofix\_case}/\code{eval\_cofix\_proj} are unrepresentable.
\end{definition}

\begin{remark}
\lead{Caveat} The module header claims rule-for-rule agreement with \code{EWcbvEval.eval} at non-block flags. Not quite: MetaRocq's \code{eval\_iota} ties the node's parameter count to the \emph{declared} \code{npars} and requires a saturated spine; \code{iota}, \code{iota\_block}, \code{iota\_sing}, \code{proj} and \code{proj\_block} all check neither. Determinism is unaffected; \WcbvEval{} is strictly \emph{more permissive} on ill-formed input.
\end{remark}

\begin{definition}[The block-form abbreviations]
\label{def:Eval}
\lean{LeanToLambdaBox.Eval, LeanToLambdaBox.EvalProp}
\leanok
\uses{def:WcbvEval, def:eraseFlags}
\lead{In short} $\mathtt{Eval}\ \Gamma = \WcbvEval\ \Gamma\ \mathtt{blockFlags}$, $\mathtt{EvalProp}\ \Gamma = \WcbvEval\ \Gamma\ \mathtt{propBlockFlags}$: block-form evaluation, propositional cases off resp.\ on, so the \code{optimize} pass's correctness statement (Chapter~\ref{chap:lowering}) can relate the two.
\end{definition}

\section{Metatheory}
\label{sec:lambdabox:metatheory}

\code{Semantics/Metatheory.lean} requires every result here to be \code{sorryAx}-free; since none of them mentions lean4lean, the measured footprints confirm it.

\begin{lemma}[Stuckness is disjoint from the spine rules]
\label{lem:isStuckApp_construct_spine}
\lean{LeanToLambdaBox.isStuckApp_construct_spine, LeanToLambdaBox.isStuckApp_fix_spine, LeanToLambdaBox.isStuckApp_construct, LeanToLambdaBox.isStuckApp_fix_bare}
\leanok
\uses{def:isStuckApp, def:LBTerm-mkApps}
\lead{In short} At any flags, an applied-form constructor spine, a bare block constructor node, and a bare fixpoint head are never stuck; under guarded recursion, neither is a fixpoint spine. Keeps congruence disjoint from the constructor and fixpoint rules.
\end{lemma}
\begin{proof}
\leanok
\uses{lem:LBTerm-mkApps_concat, def:atomValue, def:isStuckApp}
\lead{Idea} Unfold \code{isStuckApp}: the spine-head lemma reduces the head to the bare \code{construct}/\code{fix} node, so the negation collapses; bare cases by \code{simp}.
\end{proof}

\begin{theorem}[Values are final]
\label{thm:value_final}
\lean{LeanToLambdaBox.value_final}
\leanok
\uses{def:Value, def:WcbvEval}
\lead{In short} Every weak call-by-value normal form \hl{evaluates to itself}.
\begin{itemize}
\item \lead{Given} $\mathtt{Value}\ \Gamma\ fl\ v$;
\item \lead{Then} $\Gamma \vdash v \Downarrow_{fl} v$.
\end{itemize}
\lead{Lean} MetaRocq's \code{value\_final}; flag-parametric, no specialisation to \code{eraseFlags}.
\lead{Status} \stProved{}.
\end{theorem}
\begin{proof}
\leanok
\uses{def:Value, def:WcbvEval, def:atomValue}
\lead{Idea} Induction on the \code{Value} derivation, mapping each clause to its rule: \code{atom} to the atom rules; \code{construct\_block} to \code{construct}; \code{construct\_nil}/\code{construct\_app\_val} to \code{construct\_atom}/\code{construct\_app}; \code{app\_stuck}/\code{fix\_app\_val} to \code{app\_cong}/\code{fix\_stuck}, each spine clause reaching \code{mkApps} shape via its premise equation.
\end{proof}

\begin{theorem}[Evaluation produces values]
\label{thm:eval_to_value}
\lean{LeanToLambdaBox.eval_to_value}
\leanok
\uses{def:WcbvEval, def:Value}
\lead{In short} If $t$ evaluates to $v$, then \hl{$v$ is a value}: MetaRocq's \code{eval\_to\_value}, flag-parametric.
\begin{itemize}
\item \lead{Given} $\Gamma \vdash t \Downarrow_{fl} v$;
\item \lead{Then} $\mathtt{Value}\ \Gamma\ fl\ v$.
\end{itemize}
\lead{Status} \stProved{}.
\end{theorem}
\begin{proof}
\leanok
\uses{def:Value, def:WcbvEval, def:atomValue}
\lead{Idea} Induction on the evaluation: atom rules land in \code{atom}; recursive-premise rules inherit the hypothesis; \code{construct\_app}, \code{fix\_stuck}, \code{app\_cong} each build a value from a spine, feeding their side condition into the matching \code{Value} clause.
\end{proof}

\begin{theorem}[Determinism]
\label{thm:eval_deterministic}
\lean{LeanToLambdaBox.eval_deterministic}
\leanok
\uses{def:WcbvEval}
\lead{In short} A \lbox{} term has \hl{at most one value} under a fixed environment and flags.
\begin{itemize}
\item \lead{Given} $\Gamma \vdash t \Downarrow_{fl} v$ and $\Gamma \vdash t \Downarrow_{fl} v'$;
\item \lead{Then} $v = v'$.
\end{itemize}
\lead{Lean} MetaRocq's \code{eval\_deterministic}.
\lead{Status} \stProved{}.
\end{theorem}
\begin{proof}
\leanok
\uses{lem:LBTerm-mkApps_construct_inj, lem:LBTerm-mkApps_concat, lem:isStuckApp_construct_spine, def:isStuckApp, def:LBTerm-mkApps, def:WcbvEval}
\lead{Idea} Induction on the first derivation, inverting the second; the seven \code{.app}-subject rules separate by the evaluated function head's shape.
\lead{Steps}
\begin{enumerate}
\item spine identity: the spine-head lemma plus the two \code{mkApps} injectivity lemmas;
\item constructor and fixpoint spines kept apart by the disjointness lemma;
\item congruence excluded against the spine rules by the stuckness lemmas;
\item block/applied pairs and the propositional rules never collide: contradictory flag guards.
\end{enumerate}
\end{proof}

\begin{corollary}[A value evaluates only to itself]
\label{cor:eval_value}
\lean{LeanToLambdaBox.eval_value}
\leanok
\uses{def:Value, def:WcbvEval}
\lead{In short} If $v$ is a value and $v \Downarrow v'$, then \hl{$v = v'$}: MetaRocq's \code{eval\_value}.
\end{corollary}
\begin{proof}
\leanok
\uses{thm:value_final, thm:eval_deterministic}
\lead{Idea} Finality gives $v \Downarrow v$; determinism applied to that and the hypothesis gives $v = v'$.
\end{proof}

\begin{lemma}[Derivation irrelevance]
\label{lem:eval_unique}
\lean{LeanToLambdaBox.eval_unique}
\leanok
\uses{def:WcbvEval}
\lead{In short} Any two derivations of the same evaluation judgement are \hl{equal}.
\end{lemma}
\begin{proof}
\leanok
\uses{def:WcbvEval}
\lead{Idea} By \code{rfl}: \WcbvEval{} is \code{Prop}-valued, so this is proof irrelevance; MetaRocq's \code{eval\_unique} needs a real argument, its relation being \code{Type}-level.
\end{proof}

\begin{lemma}[A spine's head evaluates]
\label{lem:WcbvEval-head_value_of_mkApps}
\lean{LeanToLambdaBox.WcbvEval.head_value_of_mkApps}
\leanok
\uses{def:WcbvEval, def:LBTerm-mkApps}
\lead{In short} If a spine evaluates, so does its head: from $\mathtt{mkApps}\ f\ as \Downarrow r$, some $fv$ has $f \Downarrow fv$.
\end{lemma}
\begin{proof}
\leanok
\uses{def:WcbvEval, def:LBTerm-mkApps}
\lead{Idea} Induction on the argument list: the empty case is the hypothesis; in the cons case the spine is an \code{.app} node, and every application rule has a function-evaluation premise, so inverting yields the shorter spine's evaluation.
\end{proof}

\begin{proposition}[Applications depend on subterms only through values]
\label{prop:WcbvEval-app_congr}
\lean{LeanToLambdaBox.WcbvEval.app_congr}
\leanok
\uses{def:WcbvEval}
\lead{In short} An application's evaluation reads its subterms only through their \hl{values}, never their syntax.
\begin{itemize}
\item \lead{Given} $f,f'$ with a common value and $a,a'$ with a common value;
\item \lead{Then} if $f\,a \Downarrow r$ then $f'\,a' \Downarrow r$.
\end{itemize}
\lead{Lean} Needed by the erasure simulation (Chapter~\ref{chap:correctness}) to replace a source spine's components by related target ones.
\end{proposition}
\begin{proof}
\leanok
\uses{thm:eval_deterministic, def:WcbvEval}
\lead{Idea} Case analysis, not induction: every \code{.app} rule has one function premise and one argument premise, neither inspecting syntax. Determinism rewrites each to the common value, and the same rule re-applies to the primed terms.
\end{proof}

\begin{proposition}[The spine form of congruence]
\label{prop:WcbvEval-mkApps_congr}
\lean{LeanToLambdaBox.WcbvEval.mkApps_congr}
\leanok
\uses{def:WcbvEval, def:LBTerm-mkApps}
\lead{In short} If two argument lists are pointwise related by ``a common value'' and the two heads share a value, an evaluation of the first spine is one of the second, with the \hl{same result}.
\end{proposition}
\begin{proof}
\leanok
\uses{prop:WcbvEval-app_congr, lem:WcbvEval-head_value_of_mkApps, thm:eval_deterministic}
\lead{Idea} Induction along the pointwise list relation: the nil case is the head hypothesis; in the cons case the head-evaluation lemma supplies a value for the extended head, congruence replaces $f\,a$ by $f'\,a'$, and the hypothesis finishes the shorter lists.
\end{proof}

\begin{lemma}[Boxes absorb spines]
\label{lem:WcbvEval-mkApps_box}
\lean{LeanToLambdaBox.WcbvEval.mkApps_box}
\leanok
\uses{def:WcbvEval, def:LBTerm-mkApps}
\lead{In short} A spine whose head evaluates to $\Box$ evaluates to $\Box$, if every argument evaluates (values discarded): $\Box\,x\to\Box$ [L, Def.~5] iterated, the box arm of the erasure simulation (Chapter~\ref{chap:correctness}).
\end{lemma}
\begin{proof}
\leanok
\uses{def:WcbvEval, def:LBTerm-mkApps}
\lead{Idea} Induction on the argument list, applying the box rule once per argument.
\end{proof}

\begin{proposition}[Turning propositional cases on is a weakening]
\label{prop:WcbvEval-propcase_weaken}
\lean{LeanToLambdaBox.WcbvEval.propcase_weaken}
\leanok
\uses{def:WcbvEval, def:eraseFlags}
\lead{In short} Every evaluation at \code{eraseFlags} is one at \code{entryFlags}: no rule carries a \hl{negative} prop-case guard, so enabling the propositional rules only adds derivations.
\lead{Lean} The direction a consumer of the emitted program needs, peregrine's \code{untyped\_transform\_pipeline} declaring its input at \code{entryFlags}.
\end{proposition}
\begin{proof}
\leanok
\uses{def:WcbvEval, def:isStuckApp, def:eraseFlags}
\lead{Idea} Induction on the derivation, re-applying each rule at the new flags; arms whose guard differs are refuted, not rebuilt (block-form, \code{fix\_unguarded}, propositional arms cannot occur in an \code{eraseFlags} derivation); congruence needs \code{isStuckApp} recomputed, immediate since guarded-fix agrees.
\end{proof}

\subsection*{Non-vacuity}

\code{Semantics/Metatheory.lean} closes with witnesses against vacuous truth: a concrete evaluation $(\lambda.\,\#0)\,\Box \Downarrow \Box$, inhabitation of \code{Value} and \WcbvEval{}, and a ``fires'' witness per hypothesis-bearing lemma. \code{propcase\_weaken} and the four stuckness lemmas (Lemma~\ref{lem:isStuckApp_construct_spine}) are uncovered; two more (determinism, Corollary~\ref{cor:eval_value}) are the contentless $\Box=\Box$ and are omitted. The remaining two say something new and get nodes.

\begin{lemma}[Non-vacuity: an applied-form constructor with a parameter]
\label{lem:construct_app_fires}
\lean{LeanToLambdaBox.acKn, LeanToLambdaBox.acIid, LeanToLambdaBox.acOIB, LeanToLambdaBox.ac_arity, LeanToLambdaBox.ac_nil, LeanToLambdaBox.construct_app_fires, LeanToLambdaBox.construct_app_value}
\leanok
\uses{def:WcbvEval, def:constructorArity, def:eraseFlags, def:Value}
\lead{In short} A concrete witness at \code{eraseFlags}: one inductive, constructor \code{mk} with one parameter and one field.
\begin{itemize}
\item arity $1+1=2$ by \code{rfl}: \code{constructorArity} genuinely adds the parameter count; the bare head is a value;
\item $((\mathtt{mk})\,\Box)\,\Box \Downarrow \mathtt{mkApps}\ \mathtt{mk}\ [\Box,\Box]$ really fires, and the resulting spine is a \code{Value}.
\end{itemize}
\end{lemma}
\begin{proof}
\leanok
\uses{def:WcbvEval, thm:eval_to_value, def:constructorArity}
\lead{Idea} Arity by \code{rfl}; two \code{construct\_app} steps (arity bound by \code{decide}) on \code{construct\_atom}; the value follows from Theorem~\ref{thm:eval_to_value}.
\end{proof}

\begin{lemma}[Non-vacuity: the fixpoint substitution order]
\label{lem:mf_fixSubst}
\lean{LeanToLambdaBox.mfDefs, LeanToLambdaBox.mf_fixSubst, LeanToLambdaBox.mf_fix_stuck, LeanToLambdaBox.mf_fix_value}
\leanok
\uses{def:LBTerm-fixSubst, def:WcbvEval, def:Value}
\lead{In short} For a two-member block, $\mathtt{fixSubst}\ defs = [\mathtt{fix}\ defs\ 1, \mathtt{fix}\ defs\ 0]$ by \code{rfl} (only size $\ge 2$ pins the reversal down); applied to one argument at principal-argument index $1$, it is stuck and yields a fixpoint-spine value.
\end{lemma}
\begin{proof}
\leanok
\uses{def:LBTerm-fixSubst, def:WcbvEval, thm:eval_to_value}
\lead{Idea} \code{rfl} for the order; the stuck-fixpoint rule (empty spine, index $1$) for the evaluation; Theorem~\ref{thm:eval_to_value} for the value.
\end{proof}

\begin{theorem}[Evaluation preserves closedness]
\label{thm:WcbvEval-lbClosed}
\lean{LeanToLambdaBox.WcbvEval.lbClosed, LeanToLambdaBox.wcbvEval_lbClosed_fires_delta}
\leanok
\uses{def:WcbvEval, def:LBClosed, def:LBTerm-envLookup}
\lead{In short} The value of a closed term is \hl{closed}, provided every declared constant body is closed.
\begin{itemize}
\item \lead{Given} every declared constant body is closed, and $t$ is closed with $\Gamma \vdash t \Downarrow_{fl} v$;
\item \lead{Then} $v$ is closed.
\end{itemize}
\lead{Lean} Unavoidable: at a delta step the subject is a bare constant, closed at every bound, whose body need not be. Flag-parametric.
\lead{Status} \stProved{}.
\end{theorem}
\begin{proof}
\leanok
\uses{lem:LBClosed-substList, lem:LBClosed-shift_eq, lem:lbClosed_fix_of_bodies, def:LBClosed, def:WcbvEval, def:LBTerm-fixSubst, def:LBTerm-subst}
\lead{Idea} Induction over every arm, with an auxiliary fact first: unfolding a closed fixpoint node by $\mathtt{substList}\ (\mathtt{fixSubst}\ defs)$ yields a closed body.
\lead{Steps}
\begin{enumerate}
\item the auxiliary fact: every element of the substitution list is the same closed \code{fix} node, of the block's own length;
\item the other arms reduce to the closedness arithmetic of substitution and spines;
\item the delta arm is where the environment premise is spent.
\end{enumerate}
\end{proof}

\section{The executable evaluator}
\label{sec:lambdabox:compute}

\WcbvEval{} is \code{Prop}-valued and computes nothing; \code{lbEval} is its fuel-indexed executable twin, structurally recursive so the kernel can run it by \code{rfl}.

\begin{lemma}[Every term is its own spine]
\label{lem:LBTerm-mkApps_spine}
\lean{LeanToLambdaBox.LBTerm.mkApps_spine}
\leanok
\uses{def:LBTerm-mkApps}
\lead{In short} $\mathtt{mkApps}\ (\mathtt{spineHead}\ t)\ (\mathtt{spineArgs}\ t) = t$: the evaluator's head/arguments decomposition \hl{recovers the term}, turning a computed dispatch back into the \code{mkApps}-shaped premise the spine rules expect.
\end{lemma}
\begin{proof}
\leanok
\uses{lem:LBTerm-mkApps_concat}
\lead{Idea} Induction on the term; only \code{.app} is non-trivial, the spine-decomposition lemma turning the extended spine back into the induction hypothesis's application.
\end{proof}

\begin{definition}[The fuel-indexed executable evaluator]
\label{def:lbEval}
\lean{LeanToLambdaBox.lbEvalArgs, LeanToLambdaBox.lbEvalSpine, LeanToLambdaBox.lbEvalApp, LeanToLambdaBox.lbEvalCase, LeanToLambdaBox.lbEvalProj, LeanToLambdaBox.lbEvalStep, LeanToLambdaBox.lbEval}
\leanok
\uses{def:WcbvEval, def:constructorArity, def:isPropositionalInductive, def:isStuckApp, def:LBTerm-envLookup, def:LBTerm-subst, def:LBTerm-mkApps, def:LBTerm-fixSubst, def:WcbvFlags}
\lead{In short} An executable twin of \WcbvEval{} at the same flags, returning an optional term; fuel $0$ fails, fuel $n+1$ runs one \code{lbEvalStep} at fuel $n$.
\begin{itemize}
\item atoms return themselves; a constant's body is looked up and evaluated; \code{letIn} evaluates its value then the body; \code{construct} evaluates componentwise in block mode, nullary-and-declared only in applied mode;
\item an application evaluates both subterms, then \code{lbEvalApp}: beta, the box rule, fixpoint rules on a bare fixpoint, else \code{lbEvalSpine} --- accumulates a constructor spine under arity, unfolds a fixpoint spine at its principal-argument count, congruates else iff stuck;
\item \code{case}/\code{proj} delegate to \code{lbEvalCase}/\code{lbEvalProj}, dispatching three ways (propositional, block, spine); every unmodelled shape fails.
\end{itemize}
\lead{Lean} Structurally recursive in the fuel, so Lean's kernel reduces \code{lbEval}: a closed run is decided by \code{rfl}.
\end{definition}

\begin{remark}
\lead{Caveat} Each layer takes the recursive evaluator as an abstract parameter, keeping the fuel induction of Theorem~\ref{thm:lbEval_sound} to two lines. Three \code{DecidableEq} instances let case/projection compare identifiers --- a \emph{second} kername equality alongside \code{envLookup}'s hand-rolled one, unreconciled here.
\end{remark}

\begin{lemma}[Layered soundness of one step]
\label{lem:lbEvalStep_sound}
\lean{LeanToLambdaBox.lbEvalArgs_length, LeanToLambdaBox.lbEvalArgs_get, LeanToLambdaBox.lbEvalSpine_sound, LeanToLambdaBox.lbEvalApp_sound, LeanToLambdaBox.lbEvalCase_sound, LeanToLambdaBox.lbEvalProj_sound, LeanToLambdaBox.lbEvalStep_sound}
\leanok
\uses{def:lbEval, def:WcbvEval}
\lead{In short} Relative to a sound recursive evaluator (whatever it returns, the relation admits), each layer of the executable evaluator is \hl{sound}.
\begin{itemize}
\item argument traversal: preserves length, evaluates pointwise; spine dispatch: yields \code{construct\_app}, \code{fix\_guarded}, \code{fix\_stuck}, or \code{app\_cong};
\item application dispatch: adds beta, the box rule, the bare-fixpoint instances;
\item case/projection dispatch: yields \code{iota}/\code{iota\_block}/\code{iota\_sing}, resp.\ \code{proj}/\code{proj\_block}/\code{proj\_prop}, one whole step sound for every shape.
\end{itemize}
\end{lemma}
\begin{proof}
\leanok
\uses{lem:LBTerm-mkApps_spine, def:WcbvEval, def:lbEval, def:constructorArity, def:isPropositionalInductive, def:isStuckApp}
\lead{Idea} Each layer is a case split on the dispatch, no induction: every branch names its rule, premises from the soundness hypothesis on sub-evaluations. Lemma~\ref{lem:LBTerm-mkApps_spine} reconstitutes the \code{mkApps} shape a premise demands.
\end{proof}

\begin{theorem}[Soundness of the executable evaluator]
\label{thm:lbEval_sound}
\lean{LeanToLambdaBox.lbEval_sound}
\leanok
\uses{def:lbEval, def:WcbvEval}
\lead{In short} Whatever \code{lbEval} returns, the relation \hl{admits}; only this direction holds and only it is needed (the converse fails by fuel exhaustion).
\begin{itemize}
\item \lead{Given} $\mathtt{lbEval}\ \Gamma\ fl\ n\ t$ succeeds with $v$;
\item \lead{Then} $\Gamma \vdash t \Downarrow_{fl} v$.
\end{itemize}
\lead{Lean} How every concrete rung of the capstone (Chapter~\ref{chap:capstone}) discharges its target evaluation, by applying the theorem to a \code{rfl}.
\lead{Status} \stProved{}.
\end{theorem}
\begin{proof}
\leanok
\uses{lem:lbEvalStep_sound, def:lbEval}
\lead{Idea} Induction on the fuel: at $0$ the hypothesis says a failure equals a success (absurd); at $n+1$ it is a successful step over the fuel-$n$ evaluator, sound by the induction hypothesis, and Lemma~\ref{lem:lbEvalStep_sound} concludes.
\end{proof}

\begin{remark}
\lead{Why} Matters for Chapter~\ref{chap:trust}: \code{lbEval\_sound}'s measured footprint is exactly \code{[propext]} (no \code{sorryAx}, no \code{Classical.choice}), recorded in \code{test/ledger.expected}. The same function is the repository's differential oracle against \code{peregrine eval} on the emitted \code{.ast}.
\end{remark}

\begin{lemma}[Fuel monotonicity]
\label{lem:lbEval_mono}
\lean{LeanToLambdaBox.lbEvalArgs_mono, LeanToLambdaBox.lbEvalSpine_mono, LeanToLambdaBox.lbEvalApp_mono, LeanToLambdaBox.lbEvalCase_mono, LeanToLambdaBox.lbEvalProj_mono, LeanToLambdaBox.lbEvalStep_mono, LeanToLambdaBox.lbEval_succ, LeanToLambdaBox.lbEval_mono}
\leanok
\uses{def:lbEval}
\lead{In short} An answer found with $n$ units of fuel is found again with any \hl{larger} amount --- what makes the fuel choice of a rung, and of the differential oracle, inessential.
\end{lemma}
\begin{proof}
\leanok
\uses{def:lbEval}
\lead{Idea} Each layer transports a successful run along a pointwise-larger recursive evaluator; the successor step is step-level monotonicity at fuel $n$ versus $n+1$; the general statement is an induction on the fuel difference.
\end{proof}

\begin{remark}
\lead{Caveat} \code{Compute.lean} also carries a Peano fixture and eight examples (delta, iota, beta, fixpoint unfolding), checked by kernel reduction. No completeness theorem exists or can: the evaluator is necessarily incomplete for fuel, and order-committed --- its application arm evaluates the argument \emph{before} dispatch, sound only because every rule carries that premise.
\end{remark}

\section*{Deviations from the paper, and caveats}

\begin{itemize}
\item \textbf{Dropped side conditions.} Iota/projection rules skip MetaRocq's \code{npars}-match and saturation premises: \WcbvEval{} is more permissive than \code{EWcbvEval} on ill-formed input; determinism is unaffected.
\item \textbf{No cofixpoints, no lazy terms.} \code{LBTerm} lacks \code{tCoFix}/\code{tLazy}; MetaRocq's cofixpoint rules and \code{isLazyApp} are unrepresentable. Lean has no coinductive types, so nothing is lost.
\item \textbf{\code{kelim} and the recursivity kind are dead weight.} \code{register\_inductive} sets \code{propositional} from the source arity; \code{kelim} stays at \code{IntoAny} and the recursivity kind at \code{finite}, read by no pass.
\item \textbf{The handover is stated but unused.} Proposition~\ref{prop:WcbvEval-propcase_weaken} is what peregrine's pipeline needs; no theorem applies it.
\item \textbf{The verified perimeter stops at the \code{Program} value.} The printer is an unverified \code{partial def}: ``the emitted program satisfies \ldots'' means that value, not parsed bytes.
\item \textbf{A clean core.} \code{Semantics/} has no \code{sorry}, \code{axiom}, \code{partial}, \code{unsafe}, compiler-implemented declaration, \code{native\_decide}, or kernel-accelerated \code{decide}: witnesses are plain kernel reduction.
\end{itemize}
